\section{Literature Review}

The proposed research sits at the intersection of three distinct but interconnected bodies of scholarly literature: (1) the architecture, risk, and empirical study of DeFi lending protocols, (2) behavioral finance theory, with particular emphasis on prospect theory as a framework for understanding decision-making under risk, and (3) the emerging literature on on-chain credit risk assessment and the use of blockchain data for predictive modeling of financial outcomes. This section reviews each of these domains in turn, identifying both their individual contributions and the specific gap that motivates the proposed study.

\subsection{DeFi Lending: Architecture, Risks, and Empirical Evidence}

DeFi lending protocols represent the most capital-intensive and empirically scrutinized segment of the decentralized finance ecosystem. These protocols enable permissionless borrowing and lending of crypto-assets through smart contracts that algorithmically manage collateralization, interest rates, and liquidation events without relying on traditional financial intermediaries~\cite{schar2021decentralized2}. The fundamental mechanism is straightforward: a borrower deposits crypto-assets as collateral and may borrow another asset up to a fraction of the collateral's value, defined by the protocol's loan-to-value (LTV) parameter. If the value of the collateral declines relative to the borrowed amount—or if the value of the borrowed asset increases—the position's health factor (HF) deteriorates. Once the HF falls below the liquidation threshold (typically 1.0), any network participant can trigger a liquidation, repaying part of the debt and claiming the corresponding collateral plus a liquidation bonus~\cite{perez2020liquidations, qin2021empirical}.

The academic literature examining these mechanisms from an empirical perspective has grown considerably. \citet{perez2020liquidations} provided one of the earliest systematic analyses of liquidation events across DeFi lending platforms, documenting that liquidations occur frequently, involve substantial transaction fees, and can create cascading risk when multiple protocols are involved. \citet{qin2021empirical} extended this analysis by examining the incentive structures created by liquidation bonuses and demonstrating how these incentives can lead to competitive ``priority gas auctions'' among liquidators during periods of high market volatility. \citet{sun2022liquidity} analyzed liquidity risk in Aave, showing that large borrowing positions can create substantial utilization pressure in specific asset pools, potentially limiting the ability of smaller depositors to withdraw.

More recently, a growing body of work has focused on cross-protocol risk and the systemic dimensions of DeFi lending. \citet{tovanich2023contagion} studied financial contagion in Compound V2, constructing protocol balance sheets and demonstrating how concentrated borrowing positions can create non-trivial default risk that propagates through interconnected liquidity pools. \citet{sevim2026interoperability} extended risk modeling to multi-chain environments, showing that cross-chain interoperability introduces new categories of risk that are not captured by single-chain models. \citet{oberholzer2026toward} proposed a comprehensive nine-dimension risk assessment framework for institutional DeFi, encompassing counterparty, market, liquidity, operational, governance, legal, regulatory, smart contract, and oracle risks. \citet{iftikhar2025automated} conducted a cross-chain comparative analysis of automated risk management mechanisms in Aave and Compound, highlighting important differences in how these two dominant lending protocols handle parameter adjustments and liquidation execution.

A recurring characteristic of this literature is its focus on the protocol-level or system-level perspective. Whether analyzing liquidation cascades, interest rate dynamics, or cross-protocol dependencies, the unit of analysis is typically the protocol, the liquidity pool, or the set of interacting smart contracts. The individual borrower—as a decision-making agent whose behavior may vary systematically with risk exposure—enters the analysis primarily as a source of demand for protocol services rather than as an object of study in its own right. Understanding how borrowers actively navigate risk accumulation and whether their behavioral responses carry credit-relevant information therefore constitutes a complementary analytical perspective that the existing protocol-centric literature has not yet systematically explored.

\subsection{Prospect Theory and Financial Decision-Making Under Risk}

Since its original formulation by \citet{kahneman1979prospect} and its subsequent extension to Cumulative Prospect Theory by \citet{tversky1992advances}, prospect theory (PT) has become one of the most influential frameworks in behavioral economics for understanding how individuals make decisions under conditions of risk and uncertainty. The theory identifies several systematic deviations from the predictions of expected utility theory. First, individuals evaluate outcomes relative to a reference point, typically the status quo, rather than in terms of final wealth levels. Second, their value function is concave in the domain of gains (risk aversion) and convex in the domain of losses (loss aversion), with the disutility of losses being approximately twice as large as the utility of equivalent gains. Third, individuals exhibit diminishing sensitivity, meaning that the marginal impact of an outcome decreases as one moves farther from the reference point. Fourth, decision weights are systematically non-linear in objective probabilities, with low-probability events being overweighted and moderate-to-high-probability events being underweighted.

These theoretical predictions have received extensive empirical support across a wide range of financial domains. \citet{benartzi1995myopic} demonstrated that myopic loss aversion—the combination of loss aversion with frequent portfolio evaluation—can explain the equity premium puzzle, one of the most prominent anomalies in financial economics. \citet{barberis2013thirty}, in a comprehensive review, documented that prospect theory has been successfully applied to explain the disposition effect in stock trading, the pricing of initial public offerings, the equity premium, and a range of other financial market phenomena. The theory has become a cornerstone of behavioral finance, providing a parsimonious framework that rationalizes empirical regularities that are difficult to explain within the standard rational expectations paradigm.

However, the application of prospect theory to DeFi lending markets presents both opportunities and challenges. On the one hand, DeFi lending provides an almost ideal empirical laboratory for testing behavioral theories. The liquidation threshold—a hard-coded, protocol-defined health factor value below which positions become eligible for liquidation—serves as an objective, exogenously determined reference point that is identical across all participants. This stands in contrast to traditional financial settings, where reference points are typically unobserved, subjective, and likely to vary across individuals~\cite{barberis2013thirty}. Furthermore, the complete observability of all borrower actions in DeFi—including collateral additions, debt repayment, additional borrowing, and collateral withdrawals—means that behavioral responses to proximity to the liquidation threshold can be measured with a level of granularity that would be impossible in traditional credit markets.

On the other hand, there are significant challenges to applying prospect theory in this context. DeFi borrowers may not perceive the protocol-defined liquidation threshold as their psychological reference point; they may instead anchor on their initial entry price, a portfolio-level target return, or other benchmarks. Moreover, the rational incentive to avoid liquidation—which imposes a non-trivial financial penalty in the form of the liquidation bonus paid to the liquidator—provides a powerful non-behavioral explanation for active position management near the liquidation threshold. Disentangling rational risk management from behavioral deviations therefore requires careful empirical design, including the use of control groups, instrumental variables, or regression discontinuity approaches that exploit the quasi-exogenous variation created by the protocol-defined liquidation threshold.

\subsection{On-Chain Credit Risk Assessment}

The third body of literature relevant to this proposal concerns the use of on-chain data for credit risk assessment and default prediction. The fundamental premise of this emerging field is that the transparency of blockchain data makes it possible to construct credit risk models without relying on the centralized credit bureaus, bank account information, or traditional financial records that underpin conventional credit scoring systems.

Several recent studies have demonstrated the feasibility of this approach. \citet{ghosh2024onchain} developed a framework for on-chain credit risk scoring in DeFi, using machine learning models trained on wallet-level transaction histories, borrowing and repayment patterns, and interaction frequencies to predict the likelihood of default. Their results suggest that on-chain behavioral data contains significant credit-relevant information, even in the absence of off-chain identity information. \citet{xu2021banks} traced the evolution of lending markets from traditional banking to decentralized finance, arguing that DeFi protocols can potentially improve capital efficiency by enabling more granular and transparent risk assessment. Various survey papers have consolidated the range of DeFi risk factors that can be extracted from on-chain data, including transaction graphs, holding patterns, and interaction histories across multiple protocols~\cite{werner2022sok, jiang2023defi}.

However, the existing on-chain credit risk literature has important limitations that the current proposal seeks to address. First, the majority of studies treat credit risk as a static classification problem: given a set of features observed at time \(t\), predict whether the borrower will default at time \(t+k\). While this approach has produced useful results, it does not capture the \emph{dynamics} of how borrowers actively manage their positions as risk accumulates. A borrower who experiences a gradual decline in their health factor over several weeks and responds by systematically adding collateral and reducing leverage may face a fundamentally different risk profile than a borrower with an identical point-in-time health factor who has experienced the same decline without any adjustment activity—yet static models would assign them identical risk scores.

Second, the existing literature has focused primarily on relatively broad on-chain activity measures—such as total transaction count, account age, interaction diversity, and protocol-level aggregates—rather than on the specific behavioral processes that occur as positions approach risky thresholds. The theoretical distinction between state variables (what the position looks like at a point in time) and process variables (how the position evolved to its current state and how the borrower responded during that evolution) has not been systematically incorporated into empirical credit risk models, despite its prominence in the broader financial economics literature on the informational content of trading behavior.

\subsection{Identifying the Gap}

The preceding review identifies a clear gap at the intersection of these three bodies of literature. DeFi lending research has thoroughly documented the mechanisms of liquidation, contagion, and systemic risk from a protocol perspective but has not systematically examined borrower behavior as a source of credit signals. Behavioral finance, encapsulated by prospect theory, provides a rich theoretical framework for understanding how individuals respond to losses and reference points but has seen only limited application to DeFi lending contexts. The on-chain credit risk literature has demonstrated that blockchain data can support predictive models of borrower default but has employed predominantly static and feature-based approaches that do not capture the processual nature of borrower risk management.

The gap, stated explicitly, is: \emph{There is no systematic empirical study that examines whether borrowers' active adjustments during periods of rising position risk in DeFi lending markets contain risk information beyond that captured by conventional on-chain state variables, and whether the behavioral patterns observed are consistent with the predictions of prospect theory.}

This gap is worth addressing for both academic and practical reasons. Academically, testing whether behavioral process variables carry incremental predictive power would connect two largely separate research traditions—the protocol-centric analysis of DeFi risk and the individual-centric tradition of behavioral finance—that have the potential to inform and enrich each other. Practically, demonstrating that active borrower behavior contains independent risk signals would have immediate implications for protocol design, including the potential to complement existing risk parameters with behaviorally-informed early warning indicators.
